% Generated draft: requires mathematical and source review.
\documentclass[11pt]{article}
\usepackage[margin=25mm]{geometry}
\usepackage{fontspec}
\setmainfont{Latin Modern Roman}
\usepackage{amsmath,amssymb,mathtools,unicode-math}
\setmathfont{Latin Modern Math}
\usepackage{enumitem,needspace,xurl,hyperref}
\setlength{\parindent}{0pt}
\setlength{\parskip}{5pt}
\setlength{\emergencystretch}{3em}
\newcommand{\N}{\mathbb{N}}
\newcommand{\Z}{\mathbb{Z}}
\newcommand{\Q}{\mathbb{Q}}
\newcommand{\R}{\mathbb{R}}
\newcommand{\C}{\mathbb{C}}
\newcommand{\F}{\mathbb{F}}
\newcommand{\A}{\mathbb{A}}
\newcommand{\PP}{\mathbb{P}}
\newcommand{\E}{\mathbb{E}}
\newcommand{\eps}{\varepsilon}
\newcommand{\dd}{\,\mathrm{d}}
\begin{document}
\section*{UnsolvedMath ID: 10000052}
\textbf{Source problem code:} catalogue entry
\subsection{Definitions and mathematical statement}
Let $G$ be an infinite locally finite graph given by a triangulation of a topological open disk, with degrees bounded by a finite constant. In a standard proper triangulation model the triangular faces fill the disk and the complex is locally finite. Let $B_G(v,r)$ be its graph-distance ball.

A rooted polynomial-volume-growth convention requires constants $C,d<\infty$ and a root $o$ such that $|B_G(o,r)|\leq Cr^d$ for $r\geq1$. Requiring the same constants for every center is stronger; the source does not specify which convention it uses.

For Bernoulli site percolation at parameter $p\in[0,1]$, vertices are independently open with probability $p$. Define
\[
p_c(G)=\inf\{p:\Pr_p(o\text{ belongs to an infinite open cluster})>0\}.
\]
For a connected graph this critical threshold is independent of the chosen root. The first conjecture is that every bounded-degree disk triangulation with the intended polynomial growth satisfies
\[
p_c(G)\geq\frac12.
\]
The second question asks whether equality holds under natural sparsity assumptions on the exceptional vertices $\{v:\deg_G(v)>6\}$. No such sparsity condition is stated. Zero density, separation, summability, or uniform local constraints would be different extra hypotheses, and none is silently chosen. The second question therefore asks to identify a suitable hypothesis and prove equality, rather than asserting equality for all polynomial-growth triangulations.
\subsection{Short English statement}
Prove $p_c\geq1/2$ for bounded-degree polynomial-growth triangulations of the disk, and determine a precise sparsity condition on vertices of degree above six that forces equality.
\subsection{Sources}
\begin{itemize}
\item Local UnsolvedMath registry, permanent ID 10000052, exact source code catalogue entry, ``Critical probability of polynomial-growth disk triangulations''; attributed to Itai Benjamini, 2015. Registry status open, classification OPEN-TRIAGE.
\item Itai Benjamini, \emph{Percolation and Coarse Conformal Uniformization}, unnumbered polynomial-growth conjecture, PDF page 3, \url{https://www.wisdom.weizmann.ac.il/~itai/resistance_percolation.pdf#page=3}, preserved at \url{https://arquivo.pt/noFrame/replay/20201231041529id_/http://www.wisdom.weizmann.ac.il/~itai/resistance_percolation.pdf}. Supplied pypdf text checked against a rendered PDF, accessed July 29, 2026. No PDF/rendering was accessed here.
\item The supplied summary places the question in a Benjamini--Schramm triangulation-percolation programme and mentions Angel--Benjamini--Harel without exact references. It reports no full resolution.
\item Provenance records cross-document deduplication/source-side review and a bounded modern audit finding no complete primary result. Current status and release-reuse terms remain NEEDS\_REVIEW.
\end{itemize}
\end{document}
